% !TEX root = ../main.tex
\chapter{Local representations of a connection on the base manifold: Yang-Mills fields}\label{ch:l22}
\providecommand{\lact}{\mathbin{\triangleright}}
\providecommand{\ract}{\mathbin{\triangleleft}}
\begin{paracol}{2}

\begin{sources}
\begin{tabularx}{\linewidth}{@{}l l L@{}}
\toprule
\textbf{Video} & \textbf{Length} & \textbf{Content}\\
\midrule
\seg{22}{1}{Lecture 22: Local representations of a connection on the base manifold: Yang-Mills fields} & 1:29:33 & Recall of the connection 1-form; what ``Lie-algebra valued 1-form'' means and that pull-backs do not care about the target; a local section gives the Yang--Mills field $\omega_U$ and a local trivialisation; the Maurer--Cartan form; the frame bundle and the connection coefficients $\Gamma^i{}_{j\mu}$; the Maurer--Cartan form of $GL(d,\R)$; the gauge map between two sections and the transformation law of Yang--Mills fields; the $GL$ case and the transformation law of the $\Gamma$ symbols.\\
\bottomrule
\end{tabularx}

\smallskip
\textbf{Transcript} \file{materials/transcripts/C22-V01.txt} (YouTube automatic captions; the blackboard symbols, in particular the several Greek letters, are reconstructed from the verbal description and from \cref{ch:l21}; the captions render Maurer--Cartan as ``market on'' and ``Rialto prevail'' for ``Lie-algebra valued'').
\end{sources}

\begin{reading}[Where this lecture sits in the other notes]
\emph{Gravity and Light}, lecture 7 (connection coefficients $\Gamma^i{}_{jk}$ and their transformation law on the tangent bundle); gauge fields and the transformation law of the Yang--Mills field are not treated in the other courses (only here). It continues \cref{ch:l21} (connection, connection 1-form $\omega$ and its properties), uses the frame bundle (\cref{ch:l19}), left-invariant fields and the exponential map (\cref{ch:l13}, \cref{ch:l18}), the adjoint map (\cref{ch:l17}) and pull-backs and push-forwards (\cref{ch:l10}). Next: parallel transport (\cref{ch:l23}). (From the lecture titles; not checked in detail.)
\end{reading}

% ==================================================================
\section{Recall, and what a Lie-algebra valued 1-form is}\label{sec:ym-recall}

Last time we studied connections on a principal bundle: a connection is a distribution of horizontal subspaces, a choice (there is no canonical one), and it can be expressed conveniently as a Lie-algebra valued 1-form on the principal bundle. Today we study how this 1-form can be expressed \emph{locally} as a Lie-algebra valued 1-form on the base manifold: this is how it is usually met in physics, and depending on context it is called the local connection coefficients $\Gamma$ (general relativity) or the Yang--Mills field (elementary particle physics). The lecture is section 5.5, ``local representation of a connection''. The terminology \emph{Yang--Mills field} is used invariantly, and put in quotation marks, since it may apply to situations that a physicist would not call a Yang--Mills field in the first place \ts{22}{1}{0:07} \ts{22}{1}{0:28} \ts{22}{1}{0:49} \ts{22}{1}{1:11} \ts{22}{1}{1:32}.
\begin{recall}[Connection 1-form]
Principal $G$-bundle $(P,\pi,M)$ with right action $\ract$ and connection 1-form $\omega\colon T_pP\to T_eG$ (\cref{def:cn-omega}), satisfying (\cref{thm:cn-omega}): (a) $\omega(X^A)=A$ for the vector field $X^A$ generated by $A\in T_eG$; (b) $r_g^*\omega=(\mathrm{Ad}_{g^{-1}})_*\omega$; (c) $\omega$ smooth.
\end{recall}
\ts{22}{1}{1:53} \ts{22}{1}{2:14} \ts{22}{1}{2:34} \ts{22}{1}{2:55} \ts{22}{1}{3:16} \ts{22}{1}{3:37} \ts{22}{1}{3:58}. So much for the recall (the lecturer first writes the brackets in (b) the wrong way and corrects them: first one applies $\omega$ to $X$, the result is a Lie algebra element, and on it the push-forward of the adjoint map at $g^{-1}$ can act).

\textbf{What precisely is a Lie-algebra valued 1-form?} Somebody asked this last time \ts{22}{1}{3:58}. The simple answer: exactly what the definition says, a map from the sections of the tangent bundle, linear, not into $C^\infty(P)$, not into the real numbers or some number field, but into the Lie algebra; so one may call it a Lie-algebra valued 1-form. ``Nice, maybe suggestive name'': in what sense can one still treat it as a 1-form, and in what sense not? \ts{22}{1}{4:18} \ts{22}{1}{4:39}
\begin{answer}{In which sense is $\omega$ still a 1-form? Do pull-backs still work?}[Lie-algebra valued 1-forms]
One particular worry: all things defined for ordinary 1-forms and heavily used, like pull-backs, do they generalise? Take a map $u\colon P\to P$ such that $(u,\mathrm{id}_M)$ [?: the lecture draws a diagram with a map $f$ on the base] is a principal bundle morphism; over the same $P$ it is already an isomorphism, an \emph{automorphism} of the principal bundle (the usual two compatibility conditions, that the diagram commutes and that $u$ is equivariant, \cref{def:pb-pbmap}). If $\omega$ lives on $P$, we should be able to pull it back by $u$. But we defined the pull-back for forms whose value is in $C^\infty(P)$, or in the number field. The pull-back is, at a point,
\[
  (u^*\omega)_p(X):=\omega_{u(p)}\bigl(u_{*,p}X\bigr),
\]
defined by the push-forward of the \emph{vector} $X$, which is an ordinary vector: the pull-back does not care about the target of $\omega$, it is as Lie-algebra valued as $\omega$ is. So the one-form aspect is not harmed by changing the destination; if in doubt, go back to the definitions and check. \ts{22}{1}{4:59} \ts{22}{1}{5:19} \ts{22}{1}{5:41} \ts{22}{1}{6:02} \ts{22}{1}{6:22} \ts{22}{1}{6:43} \ts{22}{1}{7:04} \ts{22}{1}{7:24} \ts{22}{1}{7:45} \ts{22}{1}{8:06} \ts{22}{1}{8:27}
\end{answer}

% ==================================================================
\section{A local section: the Yang--Mills field and the local trivialisation}\label{sec:ym-local}

We have the bundle and the connection 1-form $\omega$, globally defined, which is no problem. In practice however we often want to restrict to part of the bundle, over a subset $U\subseteq M$ of the base manifold, for instance for calculations: $\pi\colon\pi^{-1}(U)\to U$ (the restriction is clear). Physics also takes place locally: if the base is space-time or the physical space and $P$ the principal bundle of the transformations one considers, then one does local experiments and writes down laws in a local patch; the globalisation is another issue \ts{22}{1}{8:27} \ts{22}{1}{8:48} \ts{22}{1}{9:16} \ts{22}{1}{9:40} \ts{22}{1}{10:00} \ts{22}{1}{10:20}.

So choose a \emph{local section} $\sigma\colon U\to P$, $\pi\circ\sigma=\mathrm{id}_U$ \ts{22}{1}{10:40} \ts{22}{1}{11:02}. It does two things.

\textbf{First, the Yang--Mills field} \ts{22}{1}{11:23}.
\begin{definition}[New concepts: the Yang--Mills field $\omega_U$ of a section]\label{def:ym-field}
Let $\sigma\colon U\to P$ be a local section and $\omega$ a connection 1-form on $P$. The \textbf{``Yang--Mills field''} (with respect to $\sigma$) is the Lie-algebra valued 1-form on $U\subseteq M$
\[
  \omega_U:=\sigma^*\omega ,\qquad(\omega_U)_x(v)=\omega_{\sigma(x)}\bigl(\sigma_{*,x}v\bigr)\in T_eG ,\quad v\in T_xM .
\]
\ts{22}{1}{11:46} \ts{22}{1}{12:09} \ts{22}{1}{12:30}
\end{definition}
It lives on an entirely different level than $\omega$: $\omega$ is a 1-form on the principal bundle, $\omega_U$ is a (still Lie-algebra valued) 1-form only on a local patch of the base manifold \ts{22}{1}{12:09} \ts{22}{1}{12:30} \ts{22}{1}{12:52}.

\textbf{Second, a local trivialisation} \ts{22}{1}{13:13}. Recall that a principal bundle of the form $M\times G$ is trivial, and that a principal bundle is trivial if and only if it has a global section (\cref{thm:pb-trivial-section}); in general there is none, which is why one has to restrict to a local section in the first place. But a local section gives a local trivialisation:
\begin{definition}[New concepts: the local trivialisation $h_\sigma$]\label{def:ym-triv}
The section $\sigma$ defines the map
\[
  h\colon U\times G\to P,\qquad h(m,g):=\sigma(m)\ract g ,
\]
a point $m$ of $U$ is lifted by $\sigma$ to a point of $P$, and $g$ moves one in the fibre. It is a diffeomorphism of $U\times G$ onto $\pi^{-1}(U)$, compatible with the projections and equivariant: $h(m,g\bullet g')=h(m,g)\ract g'$. The \textbf{local representation of the connection 1-form} is the pull-back $h^*\omega$, a Lie-algebra valued 1-form on $U\times G$. \ts{22}{1}{13:38} \ts{22}{1}{13:58} \ts{22}{1}{14:19} \ts{22}{1}{14:40} \ts{22}{1}{15:02} \ts{22}{1}{15:24}
\end{definition}
The lecture draws the full diagram: $U\times G\xrightarrow{\ h\ }P$, over it the projection to the first factor (not a principal bundle map, since only part of $P$ is involved, there is the right action of $G$ on the second factor with the primed action, and underneath the identity on $U$). We can pull $\omega$ back along $h$, which gives $h^*\omega$ on the local trivialisation, and along $\sigma$, which gives $\omega_U$ on the base \ts{22}{1}{15:45} \ts{22}{1}{16:06} \ts{22}{1}{16:27}. Which of the two contains all the information of $\omega$ over $U$? One would guess $h^*\omega$, and that is true; but also the Yang--Mills field $\omega_U$, living down on the base, contains the same information over $U$: both contain less than $\omega$ only because they are defined over a local patch \ts{22}{1}{16:52} \ts{22}{1}{17:13} \ts{22}{1}{17:34} \ts{22}{1}{17:55}. In physics one usually writes down only the local Yang--Mills field, in coordinates: the potential $A_\mu$ of electrodynamics, the (Lie-algebra valued) gauge potential of non-abelian gauge theories, the connection coefficients $\Gamma^\mu{}_{\nu\sigma}$ of differential geometry and general relativity: all are local shadows, local representatives, of the abstract everywhere defined connection 1-form, and they are important because they come up all the time \ts{22}{1}{17:55} \ts{22}{1}{18:15} \ts{22}{1}{18:36} \ts{22}{1}{18:57}.

\textbf{The relation} between $h^*\omega$ and $\omega_U$ \ts{22}{1}{19:19}. At a point $(m,g)$ the pull-back $h^*\omega$ takes a tangent vector of $U\times G$, that is (since $T_{(m,g)}(U\times G)=T_mU\oplus T_gG$) a pair $(v,\xi)$ of a tangent vector to the base and a tangent vector to the group, to the Lie algebra. For this one needs the Maurer--Cartan form.
\begin{coursenote}[A slip in the lecture]
While setting up the pull-back $h^*\omega$ the lecturer first writes a push-forward, then a cotangent space, then ``it takes a vector, it takes an element of $U\times G$'' and corrects himself several times (``I am sorry, it is the pull-back, I wrote down some push-forward''); what is meant is: $h_*$ maps $T_{(m,g)}(U\times G)\to T_{h(m,g)}P$, and $\omega$ eats the result. \ts{22}{1}{19:19}
\end{coursenote}
\ts{22}{1}{19:40} \ts{22}{1}{20:12} \ts{22}{1}{20:33} \ts{22}{1}{20:54}
\begin{definition}[New concepts: Maurer--Cartan form]\label{def:ym-mc}
Let $G$ be a Lie group. The \textbf{Maurer--Cartan form} is the Lie-algebra valued 1-form $\Xi$ on $G$ which at $g\in G$ is the map
\[
  \Xi_g\colon T_gG\to T_eG,\qquad\Xi_g(X):=(L_{g^{-1}})_*X ,
\]
that is: every $X\in T_gG$ is the value at $g$ of exactly one left-invariant vector field $L^A$, $L^A_g=(L_g)_*A=X$ (\cref{thm:ll-iso}), and $\Xi_g(X):=A$. Since $\Xi$ eats a tangent vector at a point of the group and yields a Lie algebra element, it is a Lie-algebra valued 1-form \emph{on the group}. \ts{22}{1}{21:15} \ts{22}{1}{21:36} \ts{22}{1}{21:56} \ts{22}{1}{22:19} \ts{22}{1}{23:22} \ts{22}{1}{23:43} \ts{22}{1}{24:03}
\end{definition}
The lecture draws the picture of the group as a sphere: at the identity $e$ the tangent space, at $g$ the tangent space; take an arbitrary vector at $g$; there is always a specific vector $A$ at the identity such that the left-invariant field generated by $A$ reproduces it at $g$. $\Xi_g$ is the inverse of this: ``you take the vector here and you find from which vector it has been produced'' \ts{22}{1}{24:24} \ts{22}{1}{24:44} \ts{22}{1}{25:05} \ts{22}{1}{25:25} \ts{22}{1}{25:47}. Here $L_g$ denotes the left multiplication $\tilde g\mapsto g\bullet\tilde g$, and its push-forward goes $T_eG\to T_gG$, as in \cref{def:ll-push-field} \ts{22}{1}{27:45} \ts{22}{1}{28:06} \ts{22}{1}{28:27}. So we now have three Lie-algebra valued 1-forms at three levels: $\omega$ on the total space of the principal bundle, $\omega_U$ on (part of) the base, and the Maurer--Cartan form on the group; but the last is fixed, a property of the group alone \ts{22}{1}{26:09} \ts{22}{1}{26:29} \ts{22}{1}{26:50}. The Maurer--Cartan form of $GL(d,\R)$ is computed in section~\ref{sec:ym-mc}; all this machinery is needed only because we want to represent the whole thing on the base space \ts{22}{1}{27:23}.
\begin{theorem}[Local form of the connection 1-form]\label{thm:ym-local}
For $(m,g)\in U\times G$, $v\in T_mU$ and $\xi\in T_gG$,
\[
  (h^*\omega)_{(m,g)}(v,\xi)=(\mathrm{Ad}_{g^{-1}})_*\bigl(\omega_U(v)\bigr)+\Xi_g(\xi) .
\]
\ts{22}{1}{21:36} \ts{22}{1}{21:56} \ts{22}{1}{22:19} \ts{22}{1}{22:41} \ts{22}{1}{23:01}
\end{theorem}
The Yang--Mills field $\omega_U$ appears and eats the tangent vector $v$, producing a Lie algebra element on which $(\mathrm{Ad}_{g^{-1}})_*$ is then applied; and an extra term which is the Maurer--Cartan form acting on $\xi$ \ts{22}{1}{22:41} \ts{22}{1}{23:01}.
\begin{proof}
\emph{(The lecture states the relation and explains the pieces, but gives no proof; the argument is added, not said in class.)} Let $A:=\Xi_g(\xi)$, so that $\xi=(L_g)_*A$. Choose a curve $c(t)=g\,\exp(tA)$ with $\dot c(0)=\xi$, and a curve $m(t)$ with $\dot m(0)=v$. Then
\[
  \begin{aligned}
  h(m(t),c(t))&=\sigma(m(t))\ract g\ract\exp(tA)\\
  \tfrac{\mathrm d}{\mathrm dt}\big|_0&=(r_g)_*\sigma_*v+X^A_{h(m,g)} .
  \end{aligned}
\]
Apply $\omega$: the first term gives $(r_g^*\omega)(\sigma_*v)=(\mathrm{Ad}_{g^{-1}})_*\omega_U(v)$ by (b); the second gives $A$ by (a). (For a general $(v,\xi)$ use linearity in $T_mU\oplus T_gG$.)
\end{proof}
(This also shows that $h^*\omega$, hence $\omega$ over $U$, is determined by $\omega_U$ alone: the two contain the same information, as said above.)

% ==================================================================
\section{Example 1: the frame bundle and the connection coefficients}\label{sec:ym-frame}

An example to motivate the choice of a section as ``localisation'' \ts{22}{1}{28:27} \ts{22}{1}{28:48} \ts{22}{1}{29:08}. Consider the frame bundle $LM$ (\cref{def:pb-frame}), where to every point all bases of the tangent space are attached. Every choice of a coordinate chart $(U,x)$ on $M$ induces a section of $LM$,
\[
  \sigma(m):=\Bigl(\frac{\partial}{\partial x^1}\Big|_m,\dots,\frac{\partial}{\partial x^{d}}\Big|_m\Bigr),\qquad d=\dim M,
\]
the frame provided by the coordinates at the point $m$ of the chart domain: choosing coordinates on the base immediately induces a section of the frame bundle \ts{22}{1}{29:28} \ts{22}{1}{29:51} \ts{22}{1}{30:12} \ts{22}{1}{30:32} \ts{22}{1}{30:53} \ts{22}{1}{31:15}. The Yang--Mills field $\omega_U=\sigma^*\omega$ is a 1-form on $U$ with values in the Lie algebra of $GL(d,\R)$, that is in the $d\times d$ matrices \ts{22}{1}{31:45} \ts{22}{1}{32:06}. In components
\[
  (\omega_U)^i{}_{j\mu}=:\Gamma^i{}_{j\mu},
\]
where $\mu=1,\dots,d$ is the 1-form index on the manifold, and $i$ (up), $j$ (down) are indices of the matrix Lie algebra of $GL$, ``the Lie algebra representation, not a $GL$ matrix'' \ts{22}{1}{32:30} \ts{22}{1}{32:51} \ts{22}{1}{33:12}. These are the connection coefficients of differential geometry (of general relativity, if one chooses an orthogonal frame, an extra restriction). All three indices run from $1$ to $d$, and very often one writes $\Gamma^\alpha{}_{\beta\gamma}$ pretending that all three have to do with the base manifold; but they do not: two of them label matrix elements of the Lie algebra, only $\mu$ is a proper 1-form index; \emph{these indices have a different logical status even if they run over the same range} \ts{22}{1}{33:32} \ts{22}{1}{33:53} \ts{22}{1}{34:13}. So this is an example of how Yang--Mills fields come up, by the choice of a section, which for the frame bundle is induced by the choice of coordinates \ts{22}{1}{34:34} \ts{22}{1}{34:55}.

% ==================================================================
\section{Example 2: the Maurer--Cartan form of $GL(d,\R)$}\label{sec:ym-mc}

The Maurer--Cartan form depends only on the group; we construct it for $GL(d,\R)$ (an important case because the Lie groups one meets are very often $GL$ or a subgroup, so that the same procedure works essentially in the same way) \ts{22}{1}{34:55} \ts{22}{1}{35:16} \ts{22}{1}{35:37} \ts{22}{1}{35:57}.

Choose coordinates on a chart domain of the group containing $g$ (the lecturer: on the connected component $GL^+$ of the identity; in fact $GL(d,\R)$ is an open subset of $\R^{d\times d}$, the determinant condition being open, so there is one chart on the whole group): $x^i{}_j\colon GL(d,\R)\to\R$, the $(i,j)$ entry of the matrix $g$ \ts{22}{1}{36:20} \ts{22}{1}{36:41} \ts{22}{1}{37:02} \ts{22}{1}{37:25} \ts{22}{1}{37:46}. A student corrects the target of the chart (it is $\R^{d\times d}$, the single coordinate functions map to $\R$) \ts{22}{1}{38:29} \ts{22}{1}{38:55}.

Consider the left-invariant vector field $L^A$ generated by $A\in T_eG$ (matrices); it exists abstractly (\cref{thm:ll-iso}), and we want to compute it. A vector field acts on functions, in particular on the coordinate functions, and by definition, at the point $g$, $L^A_g\,f=\frac{\mathrm d}{\mathrm dt}f\bigl(g\,\exp(tA)\bigr)\big|_{t=0}$: the curve $t\mapsto g\,\exp(tA)$ through $g$ is the one-parameter subgroup generated by $A$ moved by $g$ (it is the integral curve of the left-invariant field, which is how $\exp$ was defined, \cref{def:rc-exp}) \ts{22}{1}{38:08} \ts{22}{1}{39:16} \ts{22}{1}{39:36} \ts{22}{1}{39:56} \ts{22}{1}{40:17} \ts{22}{1}{40:38} \ts{22}{1}{40:59} \ts{22}{1}{41:19} \ts{22}{1}{41:40}. Then, in the matrix group, the multiplication is matrix multiplication and $\exp(tA)$ is the matrix exponential (\cref{ch:l18}), so
\[
  \begin{aligned}
  L^A_g\,x^i{}_j&=\tfrac{\mathrm d}{\mathrm dt}\big|_0\bigl(g\,e^{tA}\bigr)^i{}_j\\
  &=g^i{}_k\,A^k{}_j ,
  \end{aligned}
\]
since $g$ is a fixed matrix and $\frac{\mathrm d}{\mathrm dt}e^{tA}|_0=A$ \ts{22}{1}{42:00} \ts{22}{1}{42:21} \ts{22}{1}{42:42} \ts{22}{1}{43:02} \ts{22}{1}{43:24} \ts{22}{1}{43:46} \ts{22}{1}{44:06}. A quick check: the result is a real number depending on $i,j$, the $(i,j)$ entry of the matrix product $gA$, as it must be. Hence the components of $L^A_g$ are
\[
  L^A_g=g^i{}_k\,A^k{}_j\ \frac{\partial}{\partial x^i{}_j}\Big|_g ,
\]
all indices summed, in the basis $\partial/\partial x^i{}_j$ of $T_gG$ \ts{22}{1}{46:43} \ts{22}{1}{47:04} \ts{22}{1}{47:24} \ts{22}{1}{47:45}.
\begin{coursenote}[A moment of confusion about the definition of a vector]
At \ts{22}{1}{44:06} the lecturer doubts that the derivative can be taken ``without'' the function and says ``no no, I did not mean it''; what is meant is the definition used from the beginning: the vector tangent to a curve $\gamma$ at $p$ acts on $f$ as $(f\circ\gamma)'(0)$, with $\gamma(0)=p$, here $\gamma(0)=g$ (the brackets belong outside). With this the computation above is correct. \ts{22}{1}{44:27} \ts{22}{1}{44:48} \ts{22}{1}{45:09} \ts{22}{1}{45:42} \ts{22}{1}{46:03} \ts{22}{1}{46:23}
\end{coursenote}
\begin{proposition}[Maurer--Cartan form of $GL(d,\R)$]\label{prop:ym-mc-gl}
In the coordinates $x^i{}_j$,
\[
  \Xi_g=(g^{-1})^i{}_k\,\mathrm dx^k{}_j\ \ \text{(component $(i,j)$)}\qquad\text{at the point }g .
\]
\ts{22}{1}{48:06} \ts{22}{1}{48:27} \ts{22}{1}{48:55} \ts{22}{1}{49:15}
\end{proposition}
\ts{22}{1}{47:45} It is a form, so its components are $\mathrm dx^k{}_j$, and the result lies in the Lie algebra, with indices $(i,j)$ \ts{22}{1}{49:15}.
\begin{proof}
It suffices to check $\Xi_g(L^A_g)=A$ for all $A$, since $\Xi_g$ is defined as the inverse of $A\mapsto L^A_g$. Use the dual bases $\mathrm dx^k{}_j(\partial/\partial x^p{}_q)=\delta^k_p\delta_j^q$ and write the field with the dummy indices $p,r,q$:
\[
  \begin{aligned}
  \Xi_g(L^A_g)^i{}_j&=(g^{-1})^i{}_k\,g^p{}_rA^r{}_q\ \mathrm dx^k{}_j\Bigl(\tfrac{\partial}{\partial x^p{}_q}\Bigr)\\
  &=(g^{-1})^i{}_p\,g^p{}_r\,A^r{}_j=A^i{}_j .
  \end{aligned}
\]
\ts{22}{1}{49:37} \ts{22}{1}{49:59} \ts{22}{1}{50:20} \ts{22}{1}{50:41} \ts{22}{1}{51:02}
\end{proof}
The constants pull out, the $\mathrm dx$ are dual to the $\partial/\partial x$ by construction, giving Kronecker deltas, and the matrix multiplication $g^{-1}g$ gives the identity: ``fantastic, you plug in the left-invariant vector field generated from $A$ and you get back the $A^i{}_j$'' \ts{22}{1}{50:41} \ts{22}{1}{51:02} \ts{22}{1}{51:22}. \emph{The lesson}: to compute the Maurer--Cartan form, compute how the left-invariant field acts on chosen coordinate functions, and read off, killing the factor $g$ in front (``the $L^A$ is $gA$, so you have to get rid of the $g$ in front in order to get the $A$'') \ts{22}{1}{51:43} \ts{22}{1}{52:04}.

So, summarising \ts{22}{1}{52:04} \ts{22}{1}{52:26}: from a connection 1-form on a principal bundle, looking at part of the base manifold and choosing a section $\sigma$, we pull back $\omega$ to the Yang--Mills field $\sigma^*\omega$, a Lie-algebra valued 1-form on that domain \ts{22}{1}{52:47}.

% ==================================================================
\section{Change of section: the gauge map and the transformation of Yang--Mills fields}\label{sec:ym-gauge}

In physics we are often led to write down a Yang--Mills field because we understand the situation locally, and we can then speculate about whether a global connection 1-form can be constructed out of various local pieces: two local pieces $U_1$, $U_2$ with $U_1\cap U_2\ne\emptyset$ and a Yang--Mills field on each. There must be a compatibility condition on the overlap, and this is what we study now \ts{22}{1}{52:47} \ts{22}{1}{53:09} \ts{22}{1}{53:29} \ts{22}{1}{53:50} \ts{22}{1}{54:11}.

So let $\sigma_1,\sigma_2$ be local sections over $U_1,U_2$ (the lecturer first draws non-intersecting domains and then the intersection), $\omega$ a connection 1-form on $P$ with the two Yang--Mills fields $\omega_{U_1}=\sigma_1^*\omega$ and $\omega_{U_2}=\sigma_2^*\omega$ \ts{22}{1}{54:34} \ts{22}{1}{54:54} \ts{22}{1}{55:15}. We want the transition between the two on $U_1\cap U_2$ \ts{22}{1}{55:36}.
\begin{definition}[New concepts: the gauge map $\Omega$]\label{def:ym-gauge}
The \textbf{gauge map} (``gauge function'') is the map $\Omega\colon U_1\cap U_2\to G$ such that, for all $m\in U_1\cap U_2$, $\Omega(m)$ is the unique group element with
\[
  \sigma_2(m)=\sigma_1(m)\ract\Omega(m).
\]
\ts{22}{1}{55:56} \ts{22}{1}{56:17} \ts{22}{1}{56:38} \ts{22}{1}{56:59} \ts{22}{1}{57:20}
\end{definition}
It always exists and is uniquely defined: two points in the same fibre are related by a unique group element, because the right action is free (\cref{def:pb-free}) and the fibres are orbits \ts{22}{1}{56:59} \ts{22}{1}{57:20} (not said in class: $\Omega$ is smooth, because in a local trivialisation it is $\Omega(m)=\chi_{\sigma_1}(\sigma_2(m))$, with the map $\chi$ of \cref{thm:pb-trivial-section}). It intertwines the two choices of sections.
\begin{theorem}[Transformation of Yang--Mills fields]\label{thm:ym-gauge}
Under the above assumptions, on $U_1\cap U_2$,
\[
  \omega_{U_2}=(\mathrm{Ad}_{\Omega^{-1}})_*\,\omega_{U_1}+\Omega^*\,\Xi ,
\]
that is, for $v\in T_mM$ with $m\in U_1\cap U_2$: $(\omega_{U_2})_m(v)=(\mathrm{Ad}_{\Omega(m)^{-1}})_*\bigl(\omega_{U_1}(v)\bigr)+\Xi_{\Omega(m)}\bigl(\Omega_{*}v\bigr)$. \ts{22}{1}{57:41} \ts{22}{1}{58:03} \ts{22}{1}{58:24} \ts{22}{1}{58:47} \ts{22}{1}{59:09}
\end{theorem}
``It looks similar to the other formulas, in an entirely different formula'': it says how one connection, think of the connection coefficients, is expressed in the other frames; the first term is the adjoint action with $\Omega^{-1}$ on the first Yang--Mills field, pushed forward, the second the pull-back of the Maurer--Cartan form along $\Omega$ \ts{22}{1}{59:09} \ts{22}{1}{59:31}. If one wants one can make explicit that everything is a 1-form at the point $m$ \ts{22}{1}{59:09}.
\begin{coursenote}[The proof of \cref{thm:ym-gauge}]
The lecturer announces ``before we come to a proof of this, let us work it out again for the $GL$ case'' (\cref{sec:ym-gl}) and does not come back to it. The proof is given here, not said in class. \ts{22}{1}{59:55}
\end{coursenote}
\begin{proof}
\emph{(Added, not said in class.)} Let $v=\dot m(0)$. By the definition of $\Omega$, $\sigma_2(m(t))=\sigma_1(m(t))\ract\Omega(m(t))$, so
\[
  \sigma_{2*}v=(r_{\Omega(m)})_*\sigma_{1*}v+X^{A}_{\sigma_2(m)}
\]
with $A=\Xi_{\Omega(m)}(\Omega_*v)$ (the same computation as in the proof of \cref{thm:ym-local}). Apply $\omega$: by (b) and (a)
\[
  \omega_{U_2}(v)=(\mathrm{Ad}_{\Omega^{-1}})_*\omega_{U_1}(v)+A .
\]
\end{proof}
The extra term does not depend on $\omega$ at all, only on the group and the gauge map. Since the Maurer--Cartan form of an \emph{abelian} group is just $\mathrm d$ of the coordinates (and $\mathrm{Ad}$ is trivial), for $U(1)$ this reads $\omega_{U_2}=\omega_{U_1}+\Omega^*\Xi$, the familiar gauge transformation of the potential in electrodynamics (not said in class; the lecture mentions only that the expression ``must have been seen before with gauge transformations in physics'').

% ==================================================================
\section{The case $G=GL(d,\R)$: the transformation law of the $\Gamma$ symbols}\label{sec:ym-gl}

``Before we come to a proof, let us work it out for the $GL$ case, because that is usually the best way to understand this, since it is rather formalistic'' \ts{22}{1}{59:55} \ts{22}{1}{1:00:23}. We need two things: the pull-back $\Omega^*\Xi$ and $(\mathrm{Ad}_{\Omega^{-1}})_*$ acting on $\omega_{U_1}$, for $G=GL(d,\R)$ and the frame bundle \ts{22}{1}{1:00:43} \ts{22}{1}{1:01:04}.

\textbf{The pull-back of the Maurer--Cartan form} \ts{22}{1}{1:01:04} \ts{22}{1}{1:01:24}. $\Omega$ goes from $U_1\cap U_2$ to the group, and $\Xi$ takes $T_gG$ to $T_eG$, so the pull-back $\Omega^*\Xi$ is a Lie-algebra valued 1-form on $U_1\cap U_2$ (the lecturer first begins to say that $\Xi$ acts on $\partial/\partial x^{pq}$ ``at the point $p$'' and calls it nonsense: it acts on tangent vectors of the group, and the pull-back acts on tangent vectors of the base) \ts{22}{1}{1:01:48} \ts{22}{1}{1:02:09} \ts{22}{1}{1:02:29} \ts{22}{1}{1:02:49} \ts{22}{1}{1:03:11}. It suffices to know its action on every tangent vector $\partial/\partial x^\mu$ of the base: by the definition of the pull-back
\[
  (\Omega^*\Xi)^i{}_j\Bigl(\tfrac{\partial}{\partial x^\mu}\Bigr)=\Xi^i{}_j\Bigl(\Omega_*\tfrac{\partial}{\partial x^\mu}\Bigr)
\]
at the point $\Omega(m)$ \ts{22}{1}{1:03:32} \ts{22}{1}{1:03:53} \ts{22}{1}{1:04:15}. By \cref{prop:ym-mc-gl}, $\Xi=(\Omega^{-1})^i{}_k\,\mathrm dx^k{}_j$ at that point, and $\mathrm dx^k{}_j$ acts on the vector $\Omega_*\partial_\mu$ (a tangent vector to the group, since $\Omega$ lands in the group) by the action of the vector on the function $x^k{}_j$; by the definition of the push-forward this is $\partial_\mu(x^k{}_j\circ\Omega)$, and $x^k{}_j\circ\Omega$ is just the $(k,j)$ entry of the matrix $\Omega$ \ts{22}{1}{1:04:35} \ts{22}{1}{1:04:56} \ts{22}{1}{1:05:16} \ts{22}{1}{1:05:45} \ts{22}{1}{1:06:06} \ts{22}{1}{1:06:27} \ts{22}{1}{1:06:54} \ts{22}{1}{1:07:14} \ts{22}{1}{1:07:35} \ts{22}{1}{1:07:56} \ts{22}{1}{1:08:16}. Therefore
\[
  (\Omega^*\Xi)^i{}_{j\mu}=(\Omega^{-1})^i{}_k\ \partial_\mu\Omega^k{}_j ,
\]
\ts{22}{1}{1:08:38} \ts{22}{1}{1:08:59} \ts{22}{1}{1:09:21} \ts{22}{1}{1:09:41} \ts{22}{1}{1:10:03} and, in matrix notation (the product is matrix multiplication; $\Omega$ is a matrix, and the second factor is $\mathrm d$ of the matrix),
\[
  \Omega^*\Xi=\Omega^{-1}\,\mathrm d\Omega .
\]
``You must have seen this before with gauge transformations in physics'' \ts{22}{1}{1:10:24} \ts{22}{1}{1:10:45}. It looks as if there is no more trace of the Maurer--Cartan form, but this is the very special structure: $\Omega$ is just the gauge map interpolating between the two chosen sections \ts{22}{1}{1:11:05}. (Student question: all indices run to $d$; the lecturer: yes, $\Omega$ is a matrix, one takes the inverse, a matrix product, and the result has indices $i,j$ \ts{22}{1}{1:11:25} \ts{22}{1}{1:11:47}.)

\textbf{The adjoint term} \ts{22}{1}{1:11:47} \ts{22}{1}{1:12:08} \ts{22}{1}{1:12:31}. Recall $\mathrm{Ad}_g\colon G\to G$, $h\mapsto g\bullet h\bullet g^{-1}$, an action on the group, which fixes $e$, so that its push-forward sends $T_eG$ to $T_eG$, the Lie algebra to itself \ts{22}{1}{1:12:52} \ts{22}{1}{1:13:14} \ts{22}{1}{1:13:38}. In practice, for a matrix group,
\[
  (\mathrm{Ad}_g)_*A=g\,A\,g^{-1}
\]
with matrix multiplication on the right, a tangent vector being the derivative of a curve $t\mapsto e^{tA}$ in the group, and the push-forward of the vector being the derivative of the image curve $g\,e^{tA}g^{-1}$: the only things affected are the factors that stay in front and behind \ts{22}{1}{1:13:58} \ts{22}{1}{1:14:18} \ts{22}{1}{1:14:40} \ts{22}{1}{1:15:03} \ts{22}{1}{1:15:24} \ts{22}{1}{1:15:44}.
\begin{aside}[A warning about the star]
The lecturer insists: on the left the product is the group multiplication, on the right it is matrix multiplication; it only looks as if one could plug in another matrix, an algebra element rather than a group element, but formally $\mathrm{Ad}_g$ acts on the group and what acts on the algebra is its \emph{push-forward}. In many physics books the star is forgotten, writing $\mathrm{Ad}_gA=gAg^{-1}$, which is ``terribly misleading'', and one starts confusing the two ``only because of this numerical coincidence''. \ts{22}{1}{1:16:05} \ts{22}{1}{1:16:26}
\end{aside}
\textbf{Together} \ts{22}{1}{1:16:26} \ts{22}{1}{1:16:49}: for the frame bundle, $G=GL(d,\R)$, on $U_1\cap U_2$ the transition rule between two Yang--Mills fields on the same domain is
\[
  \omega_{U_2}=\Omega^{-1}\,\omega_{U_1}\,\Omega+\Omega^{-1}\,\mathrm d\Omega ,
\]
in components (the matrix indices $i,j$, the form index $\mu$)
\[
  \begin{aligned}
  (\omega_{U_2})^i{}_{j\mu}&=(\Omega^{-1})^i{}_k\,(\omega_{U_1})^k{}_{l\mu}\,\Omega^l{}_j\\
  &\quad+(\Omega^{-1})^i{}_k\,\partial_\mu\Omega^k{}_j .
  \end{aligned}
\]
\ts{22}{1}{1:17:12} \ts{22}{1}{1:17:32} \ts{22}{1}{1:17:52} \ts{22}{1}{1:18:13} \ts{22}{1}{1:18:34} \ts{22}{1}{1:18:54} \ts{22}{1}{1:19:15}
\begin{coursenote}[Index placement]
The captions give the placement of the $\Omega$'s ambiguously (``$\Omega^{ik}$ acting on the first Yang--Mills field with its tangent space index and its algebra index $kl$, then again $\Omega$, with respect to $\Omega^{-1}$ [and so the inverse is there]''); the version above is the one that follows from $(\mathrm{Ad}_{\Omega^{-1}})_*A=\Omega^{-1}A\,\Omega$ and from the Maurer--Cartan form of \cref{prop:ym-mc-gl} [?: reconstructed]. \ts{22}{1}{1:18:13}
\end{coursenote}
\textbf{How one uses this} \ts{22}{1}{1:19:36}. This is how one Yang--Mills field must transform into the other on the same domain. In practice one will rather have a Yang--Mills field here, here and here, and if they are to stem from the same connection 1-form on the principal bundle one needs to \emph{require} this as a transition function. This is the coordinate approach to connections: one has lots of connection coefficients, going from one to the other they must transform as above, and then there is one global connection on the principal bundle which corresponds to them \ts{22}{1}{1:19:36} \ts{22}{1}{1:19:57} \ts{22}{1}{1:20:18}. (The formula is entirely open as to which sections $\sigma_1,\sigma_2$ were chosen: it holds for any such sections, each Yang--Mills field being a pull-back of the global $\omega$.)

% ==================================================================
\section{Sections induced by coordinates: the $\Gamma$ symbols}\label{sec:ym-gamma}

The formula known for the transformation of connection coefficients is the one where the sections are induced by coordinates \ts{22}{1}{1:20:40} \ts{22}{1}{1:21:00}. So specialise: $\sigma_1\colon U_1\to LM$ and $\sigma_2\colon U_2\to LM$ the sections induced by two charts, coordinates $x^1,\dots,x^d$ on $U_1$ and $y^1,\dots,y^d$ on $U_2$, that is the frames $\partial/\partial x^i$ and $\partial/\partial y^j$ \ts{22}{1}{1:21:20} \ts{22}{1}{1:21:45} \ts{22}{1}{1:22:07}. Then the gauge map $\Omega$ is the Jacobian matrix: since $\frac{\partial}{\partial y^j}=\frac{\partial x^i}{\partial y^j}\frac{\partial}{\partial x^i}$ and $(e\ract g)_j=g^i{}_je_i$ (\cref{prop:pb-frame-principal}),
\[
  \Omega^i{}_j=\frac{\partial x^i}{\partial y^j},\qquad(\Omega^{-1})^i{}_j=\frac{\partial y^i}{\partial x^j}
\]
(the lecturer first tries the other way around, ``we need to check whether it is the right way around'', and settles on the one above) \ts{22}{1}{1:22:28} \ts{22}{1}{1:22:49} \ts{22}{1}{1:23:11}. In the transition formula above both sides still have the same $\mu$, a component with respect to the same coordinate basis: one wants to express the second Yang--Mills field also in terms of the new coordinates, so one needs yet another Jacobian, for the form index, $\partial x^\mu/\partial y^\nu$, in front of the remaining formula \ts{22}{1}{1:23:11} \ts{22}{1}{1:23:32} \ts{22}{1}{1:24:00}. The result (the lecturer checks the placement of $x$ and $y$ several times and corrects it) is
\[
  \begin{aligned}
  (\omega_{U_2})^i{}_{j\nu}&=\frac{\partial y^i}{\partial x^k}\frac{\partial x^l}{\partial y^j}\frac{\partial x^\mu}{\partial y^\nu}(\omega_{U_1})^k{}_{l\mu}\\
  &\quad+\frac{\partial y^i}{\partial x^k}\frac{\partial^2x^k}{\partial y^\nu\,\partial y^j} .
  \end{aligned}
\]
\ts{22}{1}{1:24:22} \ts{22}{1}{1:24:44} \ts{22}{1}{1:25:05} \ts{22}{1}{1:25:26} \ts{22}{1}{1:25:47} With $(\omega_U)^i{}_{j\mu}=\Gamma^i{}_{j\mu}$ this is precisely the well-known transformation formula of the $\Gamma$ symbols: each index transforms like a tensor, plus an extra (inhomogeneous) term \ts{22}{1}{1:26:07}. The apparent absence of one transformation matrix in the general formula (three indices, two Jacobians) is because there one did not assume that the $\Omega$'s have anything to do with the coordinate transformation; in the special case everything is linked and the third Jacobian is the one for the form index \ts{22}{1}{1:26:28} \ts{22}{1}{1:26:48}.

\textbf{What do we see?} \ts{22}{1}{1:26:48} \ts{22}{1}{1:27:08} The $\Gamma$ symbols of standard differential geometry (for a connection, whether one comes from a covariant derivative or not) transform in a funny, arcane way, and it is always emphasised that they do not transform as a $(1,2)$ tensor (one index up, two down, not just because of that). But this complicated object, the Yang--Mills field, comes from a very simple geometric object: a global 1-form, which is tensorial, not a 1-form on the base manifold like the Yang--Mills field but a 1-form on the principal bundle \ts{22}{1}{1:27:29} \ts{22}{1}{1:27:49}. On the principal bundle a connection is something very simple, a Lie-algebra valued 1-form with two properties; if one pushes it down to the base manifold by force, it acquires this complicated transformation behaviour. That may be where physics starts, because one looks at a patch and glues the patches obeying this rule; but then one grows up to the principal bundle and everything becomes simple in terms of the connection 1-form of \cref{ch:l21}, itself just a convenient encoding of a choice of horizontal subspaces, ``which is the actual idea of connections'' \ts{22}{1}{1:28:10} \ts{22}{1}{1:28:30}.

``In principle everything we did today we only did because we had to pay tribute to what we usually write down in physics.'' For certain calculations it may be useful, and it must be done; but next time, when we discuss parallel transport on the principal bundle and the associated bundles, and then the covariant derivative on vector bundles, we will not need this any more: we can use it, but we will go back to the nicely defined Lie-algebra valued 1-form on the principal bundle \ts{22}{1}{1:28:50} \ts{22}{1}{1:29:11}.

% ==================================================================
\section*{Key concepts}
\begin{keyconcepts}
\textbf{Retrieval questions.}
\begin{enumerate}
  \item Why does the pull-back of a Lie-algebra valued 1-form make sense? What does it not depend on?
  \item Define the Yang--Mills field $\omega_U=\sigma^*\omega$ and the local trivialisation $h(m,g)=\sigma(m)\ract g$. State the relation between $h^*\omega$ and $\omega_U$ (with a proof).
  \item What is the Maurer--Cartan form? Compute it for $GL(d,\R)$ and check it on $L^A$.
  \item Which section of the frame bundle does a chart induce? Which indices of $\Gamma^i{}_{j\mu}$ are manifold indices?
  \item Define the gauge map $\Omega$ and state the transformation law of the Yang--Mills fields (with a proof). Why does it hold for any two sections?
  \item Derive, for $GL(d,\R)$, $\Omega^*\Xi=\Omega^{-1}\mathrm d\Omega$ and $(\mathrm{Ad}_g)_*A=gAg^{-1}$; why does the star matter?
  \item Derive the transformation of the $\Gamma$ symbols from the gauge formula with $\Omega=\partial x/\partial y$.
\end{enumerate}
\textbf{Concepts.}
\begin{itemize}
  \item Yang--Mills field and local trivialisation (\cref{def:ym-field,def:ym-triv,thm:ym-local}).
  \item Maurer--Cartan form (\cref{def:ym-mc,prop:ym-mc-gl}).
  \item Gauge map and transformation law (\cref{def:ym-gauge,thm:ym-gauge}), the $GL$ case and the $\Gamma$ symbols (\cref{sec:ym-gl,sec:ym-gamma}).
\end{itemize}
\end{keyconcepts}

\flushside
\end{paracol}
